\documentclass[11pt,letterpaper]{article}
\usepackage[T1]{fontenc}
\usepackage{lmodern}
\ifdefined\pdfgentounicode\input{glyphtounicode}\pdfgentounicode=1\fi
\usepackage[hmargin=1.1in,vmargin=1in]{geometry}
\usepackage{amsmath,amssymb,booktabs,tabularx,microtype,xcolor,hyperref}
\definecolor{linkblue}{RGB}{25,65,100}
\hypersetup{colorlinks=true,linkcolor=linkblue,citecolor=linkblue,urlcolor=linkblue,
  pdftitle={Circuit complexity of small colorful patterns},
  pdfsubject={Research summary on colorful homomorphisms and ordinary weighted paths}}
\providecommand{\tightlist}{\setlength{\itemsep}{0pt}\setlength{\parskip}{0pt}}
\newcommand{\MM}{\mathsf{MM}}
\title{Circuit complexity of small colorful patterns}
\author{}
\begin{document}
\maketitle
\section{Introduction}
How many arithmetic operations does it take to count weighted copies of a fixed path?

\subsection{Colorful patterns}
Fix a simple graph \(H\), a field \(F\), and \(n\) labels in each vertex palette.
Each edge \(uv\) has an independent matrix of formal variables. The polynomial is
\[
 Z_{H,n}=\sum_{a:V(H)\to[n]} \prod_{v\in V(H)} u^v_{a(v)}
                        \prod_{uv\in E(H)} X^{uv}_{a(u),a(v)} .
\]
Write \(C_F(H,n)\) for its unrestricted binary arithmetic circuit size.
The two columns below bound the lower and upper limits of \(\log_n(1+C_F(H,n))\).
A wheel \(W_r\) has \(r\) rim vertices and one hub.

\begin{center}
\begin{tabularx}{\linewidth}{@{}Xrr@{}}
\toprule Pattern & Lower exponent & Upper exponent\\ \midrule
Forest with an edge & \(2\) & \(2\)\\
Triangle, cycle, or any cyclic partial 2-tree & \(9/4\) & \(9/4\)\\
\(K_4,\ W_4,\ K_5-e\) & \(9/4\) & \(25/8\)\\
\(K_5\) & \(9/4\) & \(4-1/300\)\\
Triangular prism and \(K_{3,3}\) & \(4/\sqrt3\) & \(25/8\)\\
\(W_5\) and \(W_6\) & \(17173/7500\) & \(25/8\)\\
Two copies of \(K_{2,3}\) joined by a $3+3$ matching & \(4/\sqrt3\) & \(25/8\)\\
Petersen graph & \(4/\sqrt3\) & \(4\)\\ \bottomrule
\end{tabularx}
\end{center}

Eliminating a degree-two vertex uses one matrix product. Conversely,
differentiating the triangle polynomial in one matrix gives the product of the other two.
This gives the exact cyclic partial 2-tree exponent. Every simple pattern with at most
seven edges and a \(K_4\) minor has cost \(\Theta_H(C_F(K_4,n))\).
Grouping the four-clique vertices into parts of sizes \(2,1,1\) gives
the upper exponent \(\omega(2,1,1)=25/8\) by the rectangular exponent formula~\cite{matrix-exponent}.
For \(K_5\), combining equal terms in a signed expansion of two label equalities
gives upper exponent \(4-1/300\). The fractional triangle-cover method with one
dense residual contraction tree has optimal costs \(25/8\) and \(4\) for these
two patterns. The five-clique circuit uses a different contraction for each of
two families of terms. Both unrestricted exponents remain open
(Appendix~\ref{five-clique-saving}). For \(K_7\) and \(K_8\), the weighted swap of Vassilevska Williams and Zamir~\cite[Section~3]{vw-zamir}
gives upper exponents \(16/3-1/1200\) and \(49/8-1/12000\), respectively.

\subsection{What three variable edges detect}

Keep a perfect matching variable in the prism or \(K_{3,3}\) and freeze the other
six edge tables. The remaining form is trilinear. Place its three inputs at
an equilateral triangle of side \(2/3\) in the Euclidean tensor
spectrum~\cite[Theorem~1.1]{euclid}. Each graph vertex contributes half the
perimeter of its three incident sites. The total is minimized when all six
frozen sites coincide at the centroid, giving lower exponent \(4/\sqrt3\).
For \(W_5\), retain one spoke and two disjoint rim edges. A rational certificate
on 23 pairs of sites gives \(17173/7500>9/4\).
With Dirac's prism-minor theorem~\cite{dirac-structure}, this leaves four 3-connected candidates for
exponent \(9/4\), namely \(K_4,W_4,K_5-e,K_5\).

On \(K_4\), keep a triangle variable. The coefficient of
\(X_{ab}Y_{bc}Z_{ca}\) is \(\sum_d U_{ad}V_{bd}W_{cd}\) for frozen spoke
matrices \(U,V,W\). Some such tensor, at fixed alphabet \(n\), requires a
matrix product of side greater than \(n^{1.077m}\) to simulate its \(m\)-fold
power by linear maps on the three inputs. For partial 2-trees every three-edge
restriction admits side \(n^{m+o(m)}\). This comparison leaves the numerical
unrestricted circuit exponent of \(K_4\) open.

\subsection{Ordinary weighted paths and matchings}

Use one symmetric table of formal edge weights on a loopless host and count
each unoriented simple path once. Curticapean--Dell--Marx expand the count
in homomorphism counts of quotients and isolate the terms using product hosts.
They apply Theorem~1.3 to paths through length six~\cite[p.~3 of the full version]{cdm}. For formal weights,
parallel edges in a quotient retain their multiplicities. The following
pointwise comparisons in characteristic zero are up to factors depending only
on the path length. \(P_i\) denotes a path on \(i\) vertices.

\begin{center}
\begin{tabularx}{\linewidth}{@{}lX@{}}
\toprule Edges & Circuit cost in characteristic zero\\ \midrule
\(1,2\) & \(\Theta(n^2)\)\\
\(3\)--\(6\) & Square matrix multiplication\\
\(7,8\) & Colorful \(K_4\)\\
\(9\) & Colorful \(K_5-e\)\\
\(10\) & Colorful \(K_5\)\\
\(11\) & Maximum of the colorful costs of \(K_6\setminus(P_4\sqcup K_2)\) and \(K_6\setminus(P_3\sqcup P_3)\)\\
\bottomrule
\end{tabularx}
\end{center}

Traversing \(K_4\) needs seven edges because its four odd vertices rule out
an Euler trail of length six. Doubling one edge suffices. Its coefficient
is \(11/2\), so characteristic eleven loses the term. At eight edges the
coefficients \(-52,-22\) restore \(K_4\) in every odd characteristic.
At nine edges characteristics three and eleven give \(K_4\) and \(W_4\), respectively.

At eleven edges, two quotients containing \(K_{3,3}\) have coefficients
\(41/2\) and \(44\). At least one survives in characteristic zero and every odd
characteristic, giving lower exponent \(4/\sqrt3\). The ordered injective
eleven-edge path with independent tables has this lower bound over every field.
Ordinary six-edge matchings have the colorful \(K_4\) cost in characteristic
zero. In characteristic two each fixed matching of at least three edges has
exponent \(9/4\).

\paragraph{Open questions.}
In characteristic zero, does the colorful four-clique have lower circuit exponent
strictly above \(9/4\)? If so, seven is the first ordinary path length to exceed
the matrix multiplication exponent. If the four-clique has lower exponent
\(9/4\), which of lengths nine, ten and eleven first exceeds it?
For ordinary paths in characteristic two, the quotient formula would require
division by two. Does the eleven-edge path still have lower exponent at least \(4/\sqrt3\)?

\clearpage
\appendix
\section{Proofs for the small-pattern bounds}\label{proofs-for-the-small-pattern-bounds}

Throughout, \(H\) is fixed, \(n\) tends to infinity, and arithmetic circuits compute formal polynomials over a field \(F\). Constants in size comparisons may depend on \(H\). Write \(C_F(H,n)\) for the binary arithmetic circuit size of the colorful polynomial with one independent edge table per edge and one unary table per vertex, each vertex having \(n\) labels. An ordinary path uses one symmetric edge table on a loopless host and counts each unoriented simple path once.

\subsection{Elimination, separators and the upper bounds}\label{elimination-separators-and-the-upper-bounds}

\textbf{Minors.} If \(J\) is a minor of \(H\), use a spanning tree in each branch set to enforce equality of its labels. Put each variable of \(J\) on one representative vertex or edge, pin vertices outside the branch sets to one label, and set unused interactions to one. Every assignment to \(J\) has exactly one extension, so \(C_F(J,n)\le C_F(H,n)\).

A separating vertex or pair also gives an upper bound. Split at a cut vertex \(u\). On each side differentiate its scalar polynomial in the unary table at \(u\) to obtain all conditional sums. Multiply the two vectors entrywise, restore the boundary unary and sum. For a separating pair \(u,v\) in a biconnected piece, add the edge \(uv\) to each side if necessary. Each augmented side is a minor because the other side supplies a \(u\)-to-\(v\) path. Differentiate in the boundary edge table to obtain the two conditional matrices, with boundary unaries set to one. Multiply entrywise, restore the original boundary weights once, and sum. By the Baur--Strassen theorem~\cite{baur-strassen}, reverse differentiation costs a constant multiple of the scalar circuit size in every characteristic. Disconnected pieces multiply. The total cost is therefore within a pattern-dependent factor of the hardest augmented piece, with \(O_H(n^2)\) additional gates when there is an edge. An independent edge table has \(n^2\) essential inputs, which already requires \(\Omega(n^2)\) binary gates.

\textbf{Forests through width three.} Removing a leaf uses one matrix-vector product, proving the forest bound. In a width-two elimination order, merging parallel factors is entrywise multiplication and eliminating a vertex with two remaining neighbors is a square matrix multiplication. Conversely, any cycle has a triangle minor. Differentiating \(\sum_{abc}X_{ab}Y_{bc}Z_{ca}\) in \(Z\) computes \(XY\). This proves the pointwise comparison with square matrix multiplication for every cyclic partial 2-tree. A simple 3-connected graph on at least five vertices needs at least eight edges by minimum degree three. Thus a graph with at most seven edges and a \(K_4\) minor has pointwise cost \(\Theta_H(C_F(K_4,n))\), by the separator reduction.

For width-three \textbf{path} decompositions we use the rectangular product in the pathwidth algorithm of Bringmann and Slusallek~\cite[Theorem~8, full version]{bringmann-slusallek}. Assign each factor to one bag. In a bag of size four the incoming and outgoing intersections each have at most three positions after removing duplicate consecutive bags. Enlarge them temporarily to triples, by broadcasting the input and summing the extra output indices. For distinct triples \((a,c,d)\) and \((b,c,d)\), assign the local pair and unary factors to \((a,c,d),(b,c,d),(a,b)\) according to their scopes. With incoming message \(M\), the update is
\[
 M'_{bcd}=R_{bcd}\sum_a X_{ab} L_{acd} M_{acd}.
\]
This is a matrix product of dimensions \(n,n,n^2\), plus \(O(n^3)\) gates. For equal triples take \(a\) to be the fourth vertex and move the incoming message to the output as a prefactor. Smaller bags use direct summation. Thus the upper exponent is \(\omega(2,1,1)=25/8\). The exact rectangular exponent~\cite{matrix-exponent}, for \(x\ge y\ge z>0\), is
\[
 \omega(x,y,z)=x+y+\frac{((x+y)z-xy)_+^2}{4xyz}.
\]
At \(z=0\) the exponent is \(x+y\).

For \(K_4\) use a single bag. For \(K_5-e\) keep the three vertices away from the missing edge and visit its endpoints in turn. For a wheel, keep the hub and one rim vertex, and sweep adjacent pairs on the remaining rim path. For \(K_{3,3}\) keep one side and visit the other. For the prism with triangles \(123,456\) and matching \(14,25,36\), use \(1234,2345,3456\). For the doubled \(K_{2,3}\), keep the left triple while visiting the two left hubs. Introduce each vertex of the right triple just before forgetting its left partner, then visit the two right hubs. Contracting all three matching edges gives \(K_{3,4}\), which contains \(K_{3,3}\). For \(K_5\) group the vertices into sets of sizes \(2,2,1\). Assign each weight to one of the three pairwise unions containing it and contract the resulting rectangular triangle. Table preparation costs \(O(n^4)\) and \(\omega(2,2,1)=4\).

For the Petersen graph use the disjointness graph of the two-subsets of \([5]\). Deleting \(12\) leaves a subdivision of \(K_{3,3}\) with branch vertices \(1i,2i\) for \(i=3,4,5\). The edges from \(1i\) to \(2j\) for \(i\ne j\) are present. The remaining three connections go through the respective complementary pair in \(\{3,4,5\}\). For the upper bound, order the vertices and then their disjoint pairs lexicographically. Root the edge contraction at edge 14, with edges numbered from zero. Its other branch is
\[
 (((((((0,2),1),(8,13)),(((6,12),9),(10,11))),(3,5)),7),4).
\]
At each join let \(s\) count vertices incident to all three edge sets (the two children and the rest) and \(x,y,z\) the vertices meeting exactly each pair. Fixing the \(s\) labels leaves a rectangular product of exponent \(\omega(x,y,z)\). Internal indices have already been summed. Substitution in the displayed tree gives join exponents \(3,25/8,4\), followed by an \(O(n^2)\) inner product. This proves the upper bound four.

\subsection{Freezing edge tables and passing to a flat}\label{freezing-edge-tables-and-passing-to-a-flat}

A pattern vertex contributes a common label \(\Delta_S(q)=\sum_{a=1}^q\bigotimes_{i\in S} e_a\) on the set \(S\) of its incident edge sites. The product \(R\) of these labels for \(q=2\) is the coefficient tensor of the edge polynomial. The \(k\)-th tensor power has alphabet \(2^k\). A fixed matrix on a frozen edge is exactly a covector at its site, allowed to couple all \(k\) binary coordinates.

A spectral point \(G\) on three outputs adds on direct sums, multiplies on tensor products, is monotone under local linear maps, and takes the value \(r\) on the full diagonal of size \(r\). Hence \(G(T)\le\operatorname{rank}(T)\). Conversely, the ordinary tensor rank of a trilinear form is at most a constant times its scalar circuit size. To see this, differentiate in the third input block and take components of bidegrees \((0,0),(1,0),(0,1),(1,1)\) in the other two. Each multiplication creates at most two new products of linear forms, and all bilinear outputs are linear combinations of these products. Thus a lower bound on the best value \(\log_{2^k}G(T)\) after freezing gives a lower circuit exponent. Supermultiplicativity gives the limit in \(k\), and padding to the preceding power of two yields the same lower limit for all \(n\).

\textbf{Contraction identity.} Put \(c(T)=\sup_\xi G(T_\xi)\), with one covector per helper site. Let \(\iota V\) insert trivial helper legs in an output tensor \(V\). Then
\[
 \lim_{k\to\infty}c(R^{\otimes k})^{1/k}
     =\min_{P:\,P(\iota V)=G(V)\ {\rm for\ all}\ V}P(R).
\]
Here \(P\) ranges over spectral points of all the sites. Restriction gives the inequality from left to right. For the reverse, let \(u\) be the limit, which exists by supermultiplicativity and the rank bound. We give the compactness argument for nonzero \(R\).

Prescribe values \(g(\iota V)=G(V)\) on any finite list of nonzero old tensors and \(g(R)=u\), and extend the prescriptions multiplicatively on formal monomials. For \(M=(\iota V)R^a\), one has \(c(M)=G(V)c(R^a)\le g(M)\) and \(\lim_k c(M^k)^{1/k}=g(M)\). Also \(c\) is monotone and additive on direct sums because covectors choose their two blocks independently. If the \(h\) prescriptions could not be met simultaneously, compactness would give a factor \(1+\eta>1\) by which some generator is too large or too small at every spectral point. Take \(M\) to be the product of the generators to a common power \(b\) with \((1+\eta)^b>8h\). Make \(2h\) variants \(M_i\) by doubling or deleting one exponent. Choose positive integers \(L,c_i\) with
\[
              \frac1{8h}<\frac{c_i g(M_i)}{L g(M)}<\frac1{4h}.
\]
At each point \(P\), some variant satisfies \(c_i P(M_i)>L P(M)\), by the discrepancy of that generator. Strict spectral comparison therefore lets a constant number of direct copies of \((\bigoplus_i c_i M_i)^k\) supply \((L M)^k\) by restriction. Integer multiples here are direct sums. Apply \(c\), expand the source, and take \(k\)-th roots. Its upper rate is \(\sum_i c_i g(M_i)< Lg(M)/2\), whereas the target has rate \(Lg(M)\). This contradiction proves finite attainment. Intersecting the closed conditions for all old tensors gives the required extension. This argument uses Strassen's spectral comparison and compactness theorems~\cite{strassen-spectrum} and the extension argument for spectral prices~\cite[Sections~2--3]{euclid}.

\textbf{Flat placement.} Use the three-output point induced by an equilateral spectral triangle of side \(2/3\). Its binary common label has log price one. By the Euclidean spectral-space theorem over \(F\) and its price calculus~\cite[Theorem~1.1 and Sections~2--5]{euclid}, binary wires have log price equal to distance, tensor prices are jointly convex along geodesics and vanish at collapsed placements, pointed rescaled ultralimits carry the limiting prices, and in a full flat a three-site binary label costs half the perimeter. We give the additional argument needed to apply the flat formula to an optimum with auxiliary sites.

In the given plane enclose the terminals in a square. Mark its corners and the two points where the horizontal through each terminal meets its boundary. Add these marks as output sites with trivial local dimensions in \(R\). Fix their joint point from the original plane. If the log price of the shared bit on the enlarged set is \(D\), normalizing divides all original log prices by \(D\). Apply the contraction identity on these outputs and multiply the minimizing extension's prices and distances by \(D\). This recovers the desired optimal rate. In its complete CAT(0) realization, flat quadrilateral rigidity fills the marked square isometrically. Equality in the side and horizontal segment distances forces the three terminals to their marked positions.

Two operations preserve the triangle while allowing a placement of no greater source price. First, near a base point \(o\), take each site \(x\) at fraction \(1/j\) along \([o,x]\), rescale prices and distances by \(j\), and take a pointed ultralimit. Convexity bounds every resulting price by its original value. A Euclidean neighborhood of \(o\) becomes a whole flat, preserving the coordinates of the tracked points originally in the square when \(o\) is interior. Each nonconstant segment from \(o\) becomes a ray, using fractions \(t/j\) at fixed time \(t\).

Second, fix a unit ray \(\rho\). From each site \(x\) take the asymptotic ray \(r_x\) to the same point at infinity. Distances between these rays at equal times stay bounded. Along them every fixed nonzero tensor price is convex and bounded above by routing its legs over wires of bounded length. It is therefore nonincreasing. In the pointed ultralimit based at \(\rho(j)\) of the unscaled spaces, \([r_x(j)]\) lies on the full parallel line \(s\mapsto[r_x(j+s)]\). Restrict to the parallel set of the limiting line of \(\rho\). The flat strip and product decomposition theorems~\cite[Theorems~2.9 and~2.11]{bridson-curvature} identify it with \(\mathbb R\times Y\), one line over each transverse point.

Apply the first operation in the square, then the second in each of the plane's coordinate directions, using the transverse factor for the second direction. The tracked plane translates rigidly and now lies in \(E\times\{y_0\}\) inside \(E\times Y\), \(E=\mathbb R^2\). To capture a helper \((v,y)\) outside this flat, apply the first operation at \((0,y_0)\). The Euclidean factor survives the ultralimit by the squared product distance. In the transverse factor \(y_0,y\) now lie on one ray. Apply the second operation along that ray, keeping the \(E\)-coordinates fixed. The old flat and this helper lie over a common transverse point in a product with Euclidean factor \(E\times\mathbb R\). Repetition captures all helpers in one flat, with the original terminal triangle and no increase in source price. For a sum of triangle perimeters we can now orthogonally project its distances to the terminal plane. This gives the lower bound needed below. Conversely, a placement of helpers inside the terminal triangle in the original flat is admissible because it preserves the convex hull and hence the shared-bit normalization.

\subsection{The prism, the five-rim wheel and the four-clique}\label{the-prism-the-five-rim-wheel-and-the-four-clique}

\textbf{A perfect matching in two cubic graphs.} Retain the prism's rungs or a perfect matching in \(K_{3,3}\). Write \(A_i\) for the three retained edge sites. Each of the other six edges joins two different matched pairs. Each unordered type \(ij\) occurs twice, and at each endpoint of \(A_i\) the incident helpers have types \(ij,ik\) with \(j\ne k\). The flat-placement argument bounds the cost below by the sum of half-perimeters over the six vertices. In the plane this is convex in the helper positions. Average over permutations of the terminal labels, accompanied by the triangle isometries, and over exchange of the prism layers or the two sides of the bipartite graph. This cannot increase the sum. The two helpers of each type coincide. Their position \(p_{ij}\) lies on the perpendicular bisector of \(A_iA_j\), at a common signed displacement \(t\) from the centroid toward that side's midpoint. With circumradius \(r=2/(3\sqrt3)\), the terminal-to-incident-helper distance is \(L(t)=\sqrt{r^2-rt+t^2}\) and the distance between distinct helper types is \(\sqrt3|t|\). The total is bounded below by
\[
                  6\min_{t\in\mathbb R}\big(L(t)+\sqrt3|t|/2\big)=6r=4/\sqrt3.
\]
Indeed \(L\) is convex with derivative \(-1/2\) at zero. The centroid placement attains equality in the original flat. This proves the two lower bounds and, by the minors above, those for the two ten-vertex patterns.

\textbf{A perfect matching in the wheel.} Number the hub \(0\) and the rim cyclically \(1,\ldots,5\). Number edge sites from zero in the order \(01,02,03,04,05,12,23,34,45,15\), and retain sites \(0,6,8\). Use the three-site price at each rim vertex. At the hub, discarding copies of its shared label bounds its price below by the average of the four triple prices on edge sites \(012,013,024,034\). In the plane each rim triple supplies weight \(1/2\) on its three pairs, and each selected hub triple supplies \(1/8\) on its pairs. Sum weights on identical pairs to get 23 capacities \(w_{ij}\).

In coordinates whose physical position is \((x,\sqrt3\,y)\), place terminals \(0,6,8\) respectively at \((1,1)/3,(0,0),(2,0)/3\). Orient the 23 pairs from smaller to larger index. The following covectors in lexicographic pair order have denominator \(10^4\). We display their numerators.
\[
\begin{array}{rrrr}
(389,4269)&(1,4321)&(-1,4321)&(-389,4269)\\
(777,8538)&(-777,8538)&(-972,1355)&(-972,1354)\\
(-1717,-3489)&(4050,5049)&(-972,-1354)&(4559,3515)\\
(-4558,3515)&(-972,-1355)&(4558,3515)&(-4559,3515)\\
(-4050,5049)&(1717,-3489)&(4050,5049)&(-4990,0)\\
(-4559,-3515)&(-4559,3515)&(4050,-5049)&
\end{array}
\]
Each pair satisfies \(3a^2+b^2\le3(10^4 w_{ij})^2\), so its covector applied to the difference of positions is at most the weighted Euclidean distance. Outgoing minus incoming numerators vanish at every helper. At terminals \(0,6,8\) they are \((0,34256),(-17218,-17128),(17218,-17128)\). Summing the inequalities gives \((34256+2\cdot17218)/30000=17173/7500\) for every placement. Contracting a rim edge yields the smaller wheel, so the bound also applies to \(W_6\). Dirac's theorem~\cite{dirac-structure}, in the form of \cite[Theorem~2]{costalonga-kingan-kingan}, lists the simple 3-connected prism-minor-free graphs as wheels, \(K_5-e,K_5\), or \(K_{3,m}\) for \(m\ge3\) with optional edges on the small side. The last family contains \(K_{3,3}\). This proves the list of four remaining candidates.

\textbf{Matrix simulation at the four-clique.} Retain a triangle. The common label at the fourth vertex gives the sum over \(d\) in the stated coefficient. In the flat there are three triangle prices with terminals \(A_i,A_j\) and helper \(p_{ij}\), and one triangle price on the helpers. Projection and the same symmetry average now give the lower bound
\[
        \min_t\big(1+3L(t)+3\sqrt3|t|/2\big)=1+2/\sqrt3 .
\]
The term one comes from the three fixed sides of length \(2/3\). The centroid again attains equality. The square matrix tensor \(M_s\) is a product of three \(s\)-state wires and has spectral value \(s^2\) on this triangle. Since \((1+2/\sqrt3)/2>1.077\), the contraction identity gives a fixed \(n=2^k\) and fixed tables yielding \(G(T)>n^{2\cdot1.077}\). A local linear restriction \(M_s\to T^{\otimes m}\) requires \(s^2\ge G(T)^m\), proving the claimed strict side bound.

For the converse, width-two elimination consumes each input once. It forms a binary tree of matrix products and entrywise products, with linear maps for broadcasting, summing or padding. Keep three input leaves and freeze the rest. At the junction of their three pairwise paths, each branch is linear in its one variable input, taking a transpose on the branch toward the old root. The resulting form is a local linear restriction of one matrix product or one diagonal tensor of size at most \(n^2\). The asymptotic subrank of \(M_n\) is \(n^2\) because powers \(M_{n^k}\) supply diagonals of size \(n^{2k-o(k)}\) by Strassen's triangle subrank construction~\cite{strassen-relative}, as described in \cite[Section~5.1]{euclid}. Increasing \(k\) to \(m+o(m)\) therefore supplies a diagonal of size \(n^{2m}\). This proves the simulation assertion for partial 2-trees.

\subsection{The weighted quotient comparison}\label{the-weighted-quotient-comparison}

Write \(h_J(X)=\sum_f\prod_{uv\in E(J)}X_{f(u)f(v)}\) for a multigraph \(J\) on a loopless symmetric weighted host, counting edges with multiplicity. For a path of \(j\) edges, let \(\pi\) partition its \(j+1\) positions and let \(J_\pi\) identify the positions in each block. Inclusion--exclusion on equalities of positions is Möbius inversion on the partition lattice. Its coefficient is \(\mu(\pi)=\prod_{B\in\pi}(-1)^{|B|-1}(|B|-1)!\). This formula follows by summing the signs of permutations with the given cycle partition. The sum over refinements of a nontrivial block vanishes. In characteristic different from two the ordinary unoriented path polynomial is
\[
             S_j(X)=\tfrac12\sum_{\pi}\mu(\pi)h_{J_\pi}(X).
\]
Loop terms vanish. For an ordinary matching in characteristic zero the same expansion divides by the pattern automorphism order \(2^m m!\). In characteristic zero all partitions giving isomorphic \(J\) have the same sign, so their total coefficient is nonzero. Each \(h_J\) is computed by the colorful circuit of its simplification, by substituting powers of the shared entries.

Following the product-host argument of Curticapean, Dell and Marx~\cite[Theorem~1.3 and Sections~2--3, full version]{cdm}, evaluate the oracle on a product host whose edge between \((i,a)\) and \((j,b)\) has weight \(X_{ij}Y_{ab}\). Each term becomes \(h_J(X)h_J(Y)\). Take a fixed-size auxiliary host with at least as many vertices as any quotient. Among isolate-free multigraphs with invertible vertex-automorphism order the \(h_J(Y)\) are linearly independent. In a putative relation take \(J\) of largest vertex count. The monomial of a labelled injective copy uses all its vertices and has coefficient \(|\operatorname{Aut}_V J|\). A different graph in the relation cannot produce it. Fixed linear combinations of formal coefficients in \(Y\) therefore isolate any surviving term. Coefficient extraction costs a fixed factor by computing over the truncated coefficient algebra, including over finite fields.

Next give each vertex of \(J\) its own palette, force all palettes to occur by inclusion--exclusion over deleted palettes, and forbid absent edge types. A surviving palette map is bijective. Mark each allowed type \(uv\) by a variable \(z_{uv}\) and extract its prescribed multiplicity \(d_{uv}\). The bijections left are the vertex automorphisms of \(J\). On this type use the factor \(1+t_{uv}X^{uv}_{ab}\) and extract degree one in every \(t_{uv}\). This gives the colorful simplification multiplied by \(|\operatorname{Aut}_V J|\prod_{uv}d_{uv}\). Unary factors may be absorbed in incident edges. Thus one obtains the reverse circuit reduction whenever those multipliers and the isolated coefficient are units. To absorb a fixed dilation of the alphabet in the pointwise comparison, partition each colorful palette into a constant number of blocks of size at most \(\lfloor n/c\rfloor\) for dilation \(c\), and sum over choices of blocks.

In odd characteristic \(p\), reduce terms whose vertex-automorphism order is divisible by \(p\). Assignments outside the fixed set of an order-\(p\) automorphism occur in equal-weight orbits of size \(p\). Fixed assignments give its vertex-orbit quotient, retaining parallel edges, or zero if there is a loop. Repeat until every automorphism group has order prime to \(p\). Each proper reduction in an isolate-free graph moves some edge, since an odd-order permutation fixing every unordered edge fixes both its endpoints. A moved edge orbit produces a loop or multiplicity at least \(p\), which cannot fall under subsequent identifications. Thus a witness with multiplicities below \(p\) and automorphism order prime to \(p\) receives no extra contributions from proper reductions.

\subsection{Why the path table stops where it does}\label{why-the-path-table-stops-where-it-does}

The quotient multigraph of a \(j\)-edge path admits an Euler trail. If a simple minor of it has \(e\) edges and \(o\) odd-degree vertices, then
\[
                        j\ \ge\ e+\max(0,(o-2)/2).
\]
Indeed the traversal starts with at most two odd vertices. Contracting an edge while retaining parallel edges never increases the odd count, whereas deleting an edge increases it by at most two. Removing an isolated vertex or a loop does not increase it. Apply the inequality to each augmented 3-connected piece from the separator argument.

\begin{itemize}
\tightlist
\item
  Up to six edges, \(K_4\) is excluded by its four odd vertices. Larger 3-connected pieces require at least eight edges. A triangle occurs at length three, and a triangle with a pendant path occurs at each subsequent length through six. This proves the square-product comparison in characteristic zero. Lengths one and two follow by forest elimination and essential inputs.
\item
  At seven and eight edges a new five-vertex piece would need eight edges, hence would be \(K_5\) minus two disjoint edges, namely \(W_4\). Its four odd vertices exclude it. At seven the only traversal with a \(K_4\) minor doubles one of the four-clique edges. Indeed the sole lost edge must be deleted to repair parity at two of its odd vertices. The word \(C,A,B,C,D,A,B,D\) realizes the traversal. The eight-edge witnesses below cover all odd characteristics.
\item
  At nine edges a six-vertex 3-connected piece would have to be cubic, with six odd vertices. On five vertices the only possibilities are \(W_4,K_5-e\). A traversal with a \(W_4\) minor must add one edge between odd rim vertices, either doubling a rim edge or adding a chord to form \(K_5-e\). At ten edges any six-vertex piece has at least nine edges. With nine it is cubic, and with ten it has at least four odd vertices. Both are excluded. A triangle, \(K_4,K_5-e,K_5\) respectively have the required lower witnesses at lengths three, seven, nine and ten, using an Euler trail after the one parity repair for \(K_4\).
\item
  At eleven edges a seven-vertex piece would have one vertex of degree four and six of degree three, violating the inequality. Eight vertices need at least twelve edges. A six-vertex piece with eleven edges has degrees \(4,4,4,4,3,3\). Its complement has degrees \(2,2,1,1,1,1\) and consists of \(P_4\sqcup K_2\) or \(P_3\sqcup P_3\). With ten edges only \(4,4,3,3,3,3\) is possible. Two degree-three vertices are nonadjacent, since otherwise the four would be an isolated clique. Add their missing edge. With nine edges the complement of the cubic piece is 2-regular and has two disjoint edges to add. Thus each six-vertex piece lies in one of the two stated eleven-edge graphs. Smaller pieces are covered by a \(K_5\) minor of \(K_6\setminus(P_3\sqcup P_3)\). View the latter as \(K_{3,3}\) plus one edge on either side and contract the cross edge between the two vertices missed by these added edges. Both eleven-edge witnesses have Euler trails and nonzero coefficients in characteristic zero. The two reductions therefore give precisely the maximum cost recorded in the table.
\end{itemize}

For the field exceptions, fix a labelled witness \(J\) with multiplicity vector \(\mathbf a\). To count vertex words using all its edges, let \(N(v,\mathbf 0)=1\) and sum \(N(w,\mathbf a-\mathbf e_{vw})\) over neighbors \(w\) with \(a_{vw}>0\). Parallel edges are indistinguishable. Sum over starts that admit an Euler trail. Each path partition giving \(J\) has \(|\operatorname{Aut}_V J|\) labelled words. The resulting coefficient is the word count times its Möbius weight divided by \(2|\operatorname{Aut}_V J|\). The weight is constant over allowed starts for each witness below. In the last two rows the sides of \(K_{3,3}\) are \(012,345\).

\begin{center}
\small
\begin{tabularx}{\linewidth}{@{}rXrrrr@{}}
\toprule
Edges & Quotient & Words & $|\mathrm{Aut}_V|$ & $\mu$ & Coefficient \\
\midrule
7 & \(K_4\), double \(AB\) & 44 & 4 & 1 & \(11/2\) \\
8 & \(K_4\), double \(AB,AC\) & 104 & 2 & \(-2\) & \(-52\) \\
8 & \(K_4\), double \(AB,CD\) & 176 & 8 & \(-2\) & \(-22\) \\
9 & \(K_5-e\) & 264 & 12 & \(-1\) & \(-11\) \\
9 & \(W_4\), double a rim edge & 120 & 2 & \(-1\) & \(-30\) \\
9 & \(K_4\), double \(AB\), attach a two-edge tail at \(C\) & 44 & 2 & 1 & 11 \\
11 & \(K_{3,3}\), add \(03,14\) & 164 & 4 & 1 & \(41/2\) \\
11 & \(K_{3,3}\), add \(01,34\) & 704 & 8 & 1 & 44 \\
\bottomrule
\end{tabularx}
\end{center}

For the seven-edge row, start at \(C\). The recurrence gives \(8,8,6\) completions after the first steps to \(A,B,D\), respectively. Starting at \(D\) gives the same total, hence \(44\) words. Independently exchanging \(A,B\) and \(C,D\) gives exactly four vertex symmetries. Each vertex occurs twice, so the four partition factors are all \(-1\). Dividing by four symmetries and two orientations gives \(44/(4\cdot2)=11/2\). To see the formal-weight reduction explicitly, first isolate this multigraph term on a product host. After the palette and edge-type extraction, the \(AB\) factor for labels \(a,b\) is \((1+tX^{AB}_{ab})^2\). Its coefficient of \(t\) is \(2X^{AB}_{ab}\). Use the other five edge tables directly. Dividing by two and the palette automorphism factor recovers the colorful four-clique whenever the isolated coefficient is nonzero and these constants are invertible.

At seven edges characteristic eleven loses its unique \(K_4\) term. A triangle with a four-edge pendant path still yields the lower bound. It has four directed vertex words, its vertex-automorphism order is two, and its single repeated vertex gives Möbius weight \(-1\). The coefficient is therefore \(-1\). At eight, \(\gcd(52,22)=2\). At nine, characteristic eleven loses \(K_5-e\) but retains \(W_4\). In characteristic three the threefold automorphism of \(K_5-e\) produces a loop and the wheel coefficient vanishes, while the tail witness gives \(K_4\). The structural upper bounds just proved then give the asserted comparisons in these fields. At eleven, \(\gcd(41,44)=1\) and the two witnesses have invertible automorphism order and multiplicities at most two in every odd characteristic. At least one yields \(K_{3,3}\), after setting the extra edges to one. This transfers \(4/\sqrt3\). The numerical threshold in characteristic zero is therefore between seven and eleven.

\subsection{Positional paths and ordinary matchings}\label{positional-paths-and-ordinary-matchings}

For an ordered injective pattern with independent positional tables let \(E_\pi\) be the unrestricted sum with positions within each block of \(\pi\) required equal. Give each block of a partition \(\rho\) a separate palette in the injective oracle. Its output \(D_\rho\) then imposes distinctness only within those blocks on the underlying labels. The same inclusion--exclusion gives \(D_\rho=\sum_{\pi\le\rho}\mu(\pi)E_\pi\). Inversion in \(\rho\) recovers \(\mu(\pi)E_\pi\), using a fixed number of calls and no division when the coefficient is a sign. To realize the ordered eleven-edge witness, use the Euler word of \(K_{3,3}\) with \(03,14\) doubled. Four vertices have degree four and two endpoints have degree three, so every vertex occurs twice. The six two-position blocks have coefficient \(+1\) in every field. Keep one table for each original edge and put one on its surplus occurrence. This proves the all-fields ordered-path lower bound.

For six ordinary matching edges in characteristic zero, every quotient has at most six edges, so the separator argument gives the \(K_4\) upper bound. Assign the six edges to the distinct edges of \(K_4\) and identify their endpoints at each clique vertex. This loopless quotient has nonzero coefficient and gives the reverse bound via the product host.

In characteristic two, first recover an edge-colored matching (one edge per table) from an ordinary \(m\)-matching. Substitute \(\sum_i t_iX^i\) as its symmetric table and extract \(\prod_i t_i\). Each bijective coloring appears once. Independent positional matrices reduce to these symmetric colored tables by adding the two oriented edge weights for each unordered host pair. Expanding one chosen edge per color gives exactly the sum over labelled endpoint orientations, including in characteristic two. Now the three distinct triangle incidences can be made from three disjoint edges by identifying their six endpoints in pairs, and the preceding palette inversion has unit coefficient. Extra matching edges can be forced on fresh, disjoint vertices with pinned tables. This gives the square-multiplication lower bound for all fixed \(m\ge3\).

For the upper bound, pad the symmetric zero-diagonal host matrix \(X\) to even order with zero weights. Put \(J_{ii}=0\) and \(J_{ij}=1\) for \(i\ne j\). Then
\[
              [t^m]\operatorname{Pf}(J+tX)=M_{m,n}(X)
              \qquad(\operatorname{char}F=2).
\]
Pfaffian signs agree, and each chosen \(m\)-matching has an odd number of perfect completions on the remaining even number of vertices (one completion for an empty remainder). Compute in \(F[X][t]/(t^{m+1})\). Split into nearly equal even blocks. Block triangular congruence of determinant one gives the product of \(\operatorname{Pf}(A)\) and \(\operatorname{Pf}(C+B^\top A^{-1}B)\) for diagonal blocks \(A,C\) and off-diagonal block \(B\). At constant order \(A_0=J_a\) satisfies \(J_a^2=I\). The constant cross block is all ones and \(\mathbf1^\top J_a\mathbf1=0\). The Schur complement again has constant term \(J\), and is alternating. Each inverse is a fixed-length geometric series about \(J_a\) since \(t^{m+1}=0\). The order-two case is a single entry. All matrix operations are therefore bounded by a fixed number of matrix products per recursive node, with \(O_m(1)\) field operations per ring operation. Summation over the recursion gives upper exponent \(9/4\).

\subsection{The four- and five-clique bounds}\label{five-clique-saving}

\textbf{Proposition.} Over every field, the upper circuit exponent of the colorful
\(K_5\) is at most \(4-1/300\).

\emph{Proof.} Absorb unaries into incident edges. Write the vertices as
\(a,b,c,d,e\), and denote each edge factor by its endpoint names.
Split the occurrences of \(a,b\) into unprimed and primed copies. The factors
\(ab,ac,bd\) use the unprimed copies, and \(a'd,b'c,a'e,b'e\) use the primed
copies. The remaining factors are \(cd,ce,de\). Insert equality on the two
copies of the ordered pair. All ten tables remain independent.

Take \(q=128\) and pad each palette to \(n=q^m\). Use
\(\ell=\lfloor m/7\rfloor\) digit positions. On one position write
\(I=P+K\), where \(P\) exchanges the two symbols and
\[
 K(a,b;a',b')=\sum_{g<h} u_{gh}(a,b)u_{gh}(a',b'),\qquad
 u_{gh}(a,b)=[(a,b)=(g,h)]-[(a,b)=(h,g)] .
\]
This is the integral identity of Jabal Ameli, Mohrin and
Nederlof~\cite[Lemma~3]{amn-six}.

\emph{Combining identical terms.}
Put \(D=\lceil\ell/12\rceil\) and \(H=\ell-D+1\), for \(\ell\ge1\).
Expand \((P+K)^{\otimes\ell}\). In terms with at least \(D\) swaps
replace every remaining \(K\) by \(I-P\). In all other terms there
are at least \(H\) copies of \(K\), and we replace every \(P\) by \(I-K\).
Collect equal tensor products before contracting any tables.
For a subset \(S\) write \(P_S\) and \(K_S\) for the tensor products
with the indicated factor on \(S\) and identity everywhere else. The result is
\[
 I^{\otimes m}=
 \sum_{\substack{S\subseteq[\ell]\\|S|\ge D}}
   (-1)^{|S|-D}\binom{|S|-1}{D-1}P_S+
 \sum_{\substack{S\subseteq[\ell]\\|S|\ge H}}
   (-1)^{|S|-H}\binom{|S|-1}{H-1}K_S .
\]
Indeed the coefficient for a fixed swap set of size \(j\) is
\(\sum_{i=D}^j(-1)^{j-i}\binom ji=(-1)^{j-D}\binom{j-1}{D-1}\).
Interchanging \(P,K\) gives the other coefficient. This identity uses
integer coefficients and holds over every field, including characteristic two.
Each subset in either sum requires only one evaluation and one scalar weight.

\emph{Evaluation of \(K_S\).} Write \(h=|S|,\ y=h/m\).
Choose the \(h\) unordered pairs in \(\bigl(\binom q2\bigr)^h\) ways. Off \(S\) the two free words \(A,B\) of length
\(m-h\) agree across their respective primed and unprimed copies.
On \(S\) one orientation word \(L\) determines both unprimed symbols and
another word \(R\) determines both primed symbols. Each has \(2^h=n^{y/7}\)
choices. The two signs depend on \(L\) and \(R\) separately.

For fixed \(A,B\), contract \(ab,ac,bd\) over \(L\) to a table on \(c,d\).
Absorb the unprimed sign and \(ab\) into one matrix. The dimensions have
exponents \((1,y/7,1)\), with cost exponent two. Multiply by \(cd\).
Preparing this table for all \(A,B\) costs exponent \(4-2y\).
The remaining factors form the tables \(b'c\, b'e\,ce\) and \(a'd\, a'e\,de\),
with the primed sign on either table. Each has at most \(n^3\) entries.
Together with the prepared table they form a rectangular triangle on
\((B,c),(A,d),(e,R)\), of exponent
\[
 \omega(2-y,2-y,1+y/7)=4-2y+C(y),\qquad
 C(y)=\frac{81y^2}{196(1+y/7)} .
\]
Including all unordered pairs gives cost exponent
\(4+C(y)+y\log_{128}(127/256)\) for this subset.

\emph{Evaluation of \(P_S\).} Put \(r=|S|/m\).
Write \(a=(A,x), b=(B,z)\) with \(x,z\) on the swapped positions.
Then \(a'=(A,z), b'=(B,x)\). First compute
\[
 M_{A,B,x,d,e}=\sum_c (ac)(b'c)(cd)(ce).
\]
The first two factors give the matrix on \((A,B,x),c\), and the last two
give the matrix on \(c,(d,e)\). Table preparation costs at most exponent
three, and the product costs
\[
 \omega(2-r,1,2)=4-r+\frac{r^2}{8(2-r)} .
\]
Multiply by \(de\). For each fixed \(A,B\), sum over \(d\) against
\((a'd)(bd)\) by multiplying on dimensions with exponents \((r,1,1+r)\).
As \(r\le1/7\), their exponent is \(2+r\) by the rectangular formula
(including \(r=0\)). There are \(n^{2-2r}\) choices of \(A,B\).
Multiply by \(ab,a'e,b'e\) and sum the remaining at most \(n^3\) entries.
The cost exponent for this subset is \(4-r+C_d(r)\), where
\(C_d(r)=r^2/(8(2-r))\).

\emph{Summing over subsets.} Write \(p=|S|/\ell\) and
\(H_2(p)=-p\log_2 p-(1-p)\log_2(1-p)\), with \(0\log 0=0\).
The number of subsets of this size is at most \(2^{\ell H_2(p)}\).
There are only \(O(m)\) sizes. Up to vanishing rounding errors, set
\(r=p/7\) or \(y=p/7\). The two excess exponents above four are bounded by
\[
\begin{aligned}
 &H_2(p)/49-r+C_d(r) &&(p\ge1/12),\\
 &H_2(p)/49+y\log_{128}(127/256)+C(y) &&(p\ge11/12).
\end{aligned}
\]
Both bounds decrease on their indicated intervals. For the first, the
derivative in \(r\) is at most
\(\log_{128}(11)-1+C_d'(r)<0\), since \(11^2<128\) and
\(C_d'(r)=r(4-r)/(8(2-r)^2)<1/8\) for \(0\le r\le1/7\).
For the second, the entropy is decreasing and \(C'(y)\le81y/98<1/7\),
whereas \(\log_{128}(127/256)<-1/7\).
At both endpoints the entropy is \(H_2(1/12)<83/200\), as verified by
\(12^{600}<2^{249}11^{550}\). Also
\(\log_{128}(127/256)\le-129/896\) by \(\ln(1-x)\le-x\) and \(\ln 2<1\).
The two respective upper bounds on the excess are therefore
\[
 \frac{83}{9800}-\frac1{84}+\frac1{8\cdot84\cdot167}
 \quad\text{and}\quad
 \frac{83}{9800}-\frac{11}{84}\frac{129}{896}+C(11/84).
\]
Each is less than \(-1/300\). For uniform rectangular estimates, round
dimension exponents up on a fixed mesh and use the finitely many padded
algorithms. Continuity then allows arbitrarily small slack. Padding the
original palette to a power of \(q\) increases it by a constant factor.
This proves the bound. \hfill\(\square\)

\paragraph{What the fractional triangle-cover program certifies.}
Brand et al.~\cite[Theorem~5 and Lemma~39]{beyond-bilinear} pay asymptotic
rank for fractionally covered wires and handle the residual graph by a
contraction tree. Their graph tensor puts an input at each vertex and a wire
on each edge. The catalogue tensor has an input at each edge and one shared
label across the inputs incident to each vertex. The following cut bound
also applies after realizing those shared labels by local restrictions.

Suppose the edge coefficient tensor is a local restriction of a tensor
product of a resource of asymptotic rank exponent \(u\) and a residual
wire network charged exponent \(v\) by one dense contraction tree.
At any trivalent node of that tree group the input sites into its three
branches. Regrouping preserves the restriction and can only decrease rank.
The residual's grouped tensor consists of the wires crossing between the
three branches, up to nonzero internal factors on individual branches.
Its asymptotic rank exponent is at most its dense join charge. More
generally a label crossing all three branches supplies a common diagonal,
one independent product per label. Thus the grouped target at each node
has asymptotic rank exponent at most \(u+v\).

For a clique the subtrees spanned by the incidences of its individual
vertices meet pairwise at edge leaves, hence have a common node by the
Helly property of subtrees. For \(K_4,K_5\) this node is internal.
If \(s\) vertex labels occur on all three branches and \(x,y,z\) on the
three different pairs of branches, the grouped target there is a direct sum
of \(n^s\) identical rectangular multiplication tensors. Its asymptotic
rank exponent is \(s+\omega(x,y,z)\) by additivity of spectral points
and Strassen's rank characterization~\cite{strassen-spectrum}. All four
or five labels have been counted. Assigning the \(s\) common labels to
the rectangle dimensions can increase \(\omega\) by at most \(s\) by
blocking. Symmetry and convexity of \(\omega\) then minimize the cost
among integer triples of this sum at the balanced triple. Consequently a
single rank charge plus one residual dense tree costs at least \(25/8\)
for \(K_4\) and \(4\) for \(K_5\). For a fractional cover this is a dual certificate. The three-site spectral
point maximizing the grouped target assigns each available fractional
triangle a log price at most its rank budget and the residual at most its
join budget. At the balanced joins the supporting wire prices are
\((15/16,5/8,5/8)\) and \((1,1,0)\), respectively, in decreasing dimension
order. These are the gradients of the rectangular exponent at
\((2,1,1)\) and \((2,2,1)\). Their dot products with those dimension triples
are \(25/8\) and \(4\).
The ordinary rectangular trees attain
these benchmarks when rectangular joins are allowed. In the proposition
the equality expansion produces a sum and the circuit chooses a tree
separately for each summand. The unrestricted lower bound for both cliques
remains the triangle-minor bound.

\paragraph{A ceiling for the three-edge test in a Euclidean flat.}
For either clique, every choice of three variable edge tables and every
normalized Euclidean terminal triangle have contracted log price at most
\(\sqrt5<9/4\), in units of \(\log_2 n\). This bounds this family of
spectral witnesses for the unrestricted lower exponent.

Indeed, among three edges on at most five vertices two meet. Write \(A,B\)
for their sites and \(C\) for the third. The terminal triangle has perimeter
two, since a Euclidean three-site shared bit costs half its perimeter.
Place all frozen sites at a common point \(p\) of the terminal hull.
The common endpoint of the first two edges costs at most the half-perimeter
of \(A,B,p\), with any further incidences charged individually. The other
four terminal incidences can be supplied by wires from \(p\). If endpoints
also coincide, enforce equality of their copies at \(p\) and discard unused
legs. Vertices with only frozen incidences cost zero at this placement.
The source price at this extension is therefore at most
\[
 \tfrac12|A-B|+\tfrac32(|A-p|+|B-p|)+2|C-p| .
\]
The contraction identity bounds the best frozen price by the cost at any
such extension.

Here is a uniform bound on the minimum over \(p\). Minimize the last three
weighted distances over the plane. Projection allows a minimizer in the
terminal hull. The optimality condition for a sum of norms gives vectors
\(u_A,u_B,u_C\) of lengths at most \(3/2,3/2,2\) with sum zero and with
the weighted distance sum equal to \(u_A\cdot A+u_B\cdot B+u_C\cdot C\).
This also holds at a coincident site by using a subgradient in its closed
ball. Set \(b=\sqrt5/2\) and \(h=b-1/2\). There is a vector \(z\) with
\[
 |z|\le h,\qquad |u_A-z|\le b,\qquad |u_B+z|\le b .
\]
To see this, map the reference points \((0,0),(b,1),(b,-1)\) to
\(0,u_A,-u_B\). All pairwise distances weakly decrease since the reference
lengths are \(3/2,3/2,2\). The point \((h,0)\) has reference distances
\(h,b,b\). For completeness, write \(v_i,w_i\) for the reference and
image points and \(r_i=|v_i-(h,0)|\). Minimize
\(\max_i(|z-w_i|^2-r_i^2)\). Its minimizer is a convex combination of active
centers with weights \(\lambda_i\). The minimum equals
\(\frac12\sum_{ij}\lambda_i\lambda_j|w_i-w_j|^2-\sum_i\lambda_i r_i^2\).
Replacing images by reference points increases this to at most zero by the
variance identity. This proves the three ball constraints.

Now send vector \(z\) along \(AB\), \(u_A-z\) along \(AC\), and \(u_B+z\)
along \(BC\). Taking scalar products with the corresponding differences
of terminal positions bounds the weighted sum by
\(h|A-B|+b(|A-C|+|B-C|)\).
Restoring the half-edge term gives \(b\) times the perimeter, hence
\(\sqrt5\). The matrix-multiplication lower bound uses a spectral point
outside this Euclidean family.

\paragraph{Nested scalar clique queries at three Euclidean terminals.}
Let \(G\) be a normalized three-terminal Euclidean spectral point. Suppose a
division-free circuit for a trilinear form \(T(X_1,X_2,X_3)\) uses \(s\) ordinary
binary gates and scalar colorful \(K_4\) or \(K_5\) gates of palette sizes
\(n_g\). Queries may be nested and intermediate values reused. Then
\[
                     G(T)\ \le\ O\left(s+\sum_g n_g^{\sqrt5}\right).
\]
In particular, allowing arithmetic around the queries preserves the preceding
Euclidean bound per query, up to the additive cost of the ordinary gates.

\emph{Proof.} First absorb unary inputs into edge tables using \(O(n_g^2)\)
binary gates per query. Substitute \(\epsilon_i X_i\) for each input block and
calculate modulo \(\epsilon_1^2,\epsilon_2^2,\epsilon_3^2\). The constant
component of each wire is a field scalar, and its singleton components are
linear forms. For a gate \(g\) with \(m_g\) scalar inputs, let \(B_{g,ij}\)
denote the part of its \(\epsilon_i\epsilon_j\) coefficient formed from two
singleton components. Its matrix rank is at most \(m_g\), since it is a
bilinear form on two lists of \(m_g\) linear forms. The rest of this
coefficient is a scalar linear combination of the input pair components.
Inductively, each pair component is a scalar linear combination of earlier
\(B_{g,ij}\), including the gate's own term.

In the coefficient of \(\epsilon_1\epsilon_2\epsilon_3\), the gate propagates
input triple components with scalar coefficients. New terms use either an
input pair and an input singleton, or three singletons. Write \(T_g^0\) for
the latter contribution. Collecting all linear multipliers of each
\(B_{g,ij}\) at the output gives
\[
 T=\sum_g c_gT_g^0+
   \sum_g\sum_{i<j} B_{g,ij} L_{g,k},\qquad \{i,j,k\}=\{1,2,3\},
\]
with \(L_{g,k}\) linear in \(X_k\). This coefficient argument also holds in
small characteristic. Each product \(B_{g,ij}L_{g,k}\) has tensor rank at
most \(m_g\). An ordinary gate has \(m_g\le2\) and \(T_g^0=0\).
For a clique query \(m_g=O(n_g^2)\). Its \(T_g^0\) is a sum over ordered
triples of distinct edge tables, since the query is multilinear in its
tables. Each of these constantly many summands is a local linear image of
a three-edge test with all other tables frozen at their constant components.
It has \(G\)-value at most \(n_g^{\sqrt5}\) by the preceding paragraph.
A spectral point is bounded by rank and subadditive under ordinary sums,
since a sum is a restriction of the direct sum. Summing the bounds proves
the claim. \hfill\(\square\)

\paragraph{A four-clique fold and its residual rank.}
A tempting modification keeps \(ab,ad\) on the unprimed side, puts
\(a'c,b'c,b'd\) on the primed side, and keeps \(cd\).
On two \(q\)-ary words of length \(m\), let \(P_S\) swap just the positions
in \(S\), leaving the other positions fixed. For \(|S|=rm\), the direct
rectangle has dimensions \((2-r,1,1)\). Indeed, with the notation
\(a=(A,x), b=(B,z), a'=(A,z), b'=(B,x)\), the factors at \(d\) use
only \((A,B,x),d\). Sum over \(z\) in the remaining factors at \((A,B,x),c\)
and close the rectangle against \(cd\). Its dense multiplication charge
has exponent
\[
 3-r+\frac{1}{4(2-r)}\ \ge\ \frac{25}{8}-\frac{15r}{16}.
\]

Suppose we approximate \(I\) by \(D=\sum_i c_iP_{S_i}\) and try to compress
the residual \(R=I-D\) in the central rectangle. For fixed \(b'\), treat
\(R(a,b;a',b')\) as a matrix with rows \((a,b)\) and columns \(a'\).
The right-side factors at \(c\) use \(a',b'\), and the extra right-side
factor at \(d\) retains \(b'\). Separate rank factorizations for each \(b'\)
would give a rectangle dimension equal to the sum of these slice ranks.
Write \(L=\sum_i q^{-|S_i|}\). Restrict the rows to \(b=b'\). In summand
\(i\) a nonzero row must satisfy \(a=b'\) on \(S_i\), leaving at most
\(nq^{-|S_i|}\) rows. On all uncovered rows \(R\) agrees with the identity
in \(a,a'\). The summed slice rank is therefore at least \(n^2(1-L)\).
If \(L<1/2\), this dimension still has exponent two. If \(L\ge1/2\),
charging each direct rectangle gives total benchmark at least
\(n^{25/8}\sum_i q^{-15|S_i|/16}\ge n^{25/8}/2\).
This covers arbitrary scalar coefficients and choices of swap sets for
this split and these two charges.

There is also a joint rank test at the same middle index. Specialize the
left table \(X^{ab}\) to all ones. For any two-wire operator \(Q\), the
matrix between the \(d\)-side index \((a,b')\) and the \(c\)-side index
\((a',t)\) is
\[
 M_Q[(a,b'),(a',t)]=[b'=t]\sum_b Q(a,b;a',b').
\]
For \(Q=I\) this is the identity of order \(n^2\).
For \(P_S\), a nonzero row requires \(a=b'\) on \(S\).
If the union of direct row supports has size \(w\), \(M_R\) therefore
contains an identity submatrix of order \(n^2-w\). Any exact joint
factorization of this middle matrix has rank at least \(n^2-w\), even if
it mixes different \(b'\) slices. The independent \(ad,b'd\) and \(a'c,b'c\)
tables can select any row and column by pinning one \(c,d\) pair, so the
specialization tests every entry. This leaves open algorithms that share
arithmetic across summands or use a different middle index.

\subsection{The seven- and eight-clique savings}\label{the-two-clique-savings}

For completeness, the two clique bounds use one folded wedge from the weighted swap of Vassilevska Williams and Zamir~\cite[Lemmas~3.1--3.3]{vw-zamir}. Write \(k=7,8\), \(h=k-2\), and \(n=2^\ell\), using padding for other sizes. A noncomplete pattern on \(k\) vertices has two nonadjacent vertices. Given the other \(h\) labels, each of their sums is a table of the form \(\sum_i\prod_j A^j_{i,a_j}\). Split the \(h\) labels into groups of sizes \(u,v\ge2\). One rectangular product of dimensions \(n^u,n,n^v\) forms this table at exponent \(u+v=h\), since \(uv\ge u+v\). Its two matrices, the second vertex's table, and all remaining factors fit this bound.

Group the clique vertices into sizes \((3,2,2)\) or \((3,3,2)\) and put the wedge \(x,z,y\) in the first group. Assign \(xz\) to one of the two matrix factors using that group and \(yz\) to the other. The common pair index \((x,y)\) is glued by an identity of order \(n^2\). A partial swap interchanges its two strings off a set of pinned digit positions. Then the occurrence of \(x\) on one side and the occurrence of \(y\) on the other are functions of a single free string. Thus the two wedge weights merge into one pair table with \(z\), and every other weight still uses at most two free strings. Each partial-swap term costs at most exponent \(h\) by the noncomplete case, including \(O_k(n^2)\) to prepare its weights.

Let \(P\) swap two bits, \(K=I-P\), and \(D=\lfloor\alpha\ell\rfloor\). Each coefficient of \((P+tK)^{\otimes\ell}\) through degree \(D\) is a sum of partial swaps, at most \(\sum_{j\le D}\binom{\ell}{j}4^j=n^{H_2(\alpha)+2\alpha+o(1)}\) in total since \(K\) has four nonzero entries. Here \(H_2\) is binary entropy. A linear combination of these coefficient operators or its transpose can be applied in \(O(n^2\operatorname{poly}(\ell))\) gates by applying one digit at a time and storing coefficients through \(D\).

When \(2\ne0\), \(P\) has three \(+1\) basis vectors and one \(-1\) basis vector. On a word with \(j\) minus signs the degree-\(d\) coefficient has eigenvalue \((-1)^j(-2)^d\binom jd\). Choose an integer interval \([b,b+D]\) centered within one of \(\ell/4\). On this interval \(Q(j)=(-1)^b\sum_{d=0}^D(-2)^d\binom{j-b}{d}=(-1)^j\). Vandermonde's identity expands \(Q\) in the functions \((-2)^d\binom jd\) with integer coefficients, also for a negative shift \(-b\). Thus a combination \(S\) of the low-degree operators with eigenvalue \((-1)^jQ(j)\) equals the identity on this interval in every odd characteristic and in characteristic zero. Its remainder \(I-S\) factors through the exceptional rows in the product eigenbasis, at most \(4^\ell\exp(-\alpha^2\ell/2+o(\ell))\) by the binomial tail for \(\operatorname{Bin}(\ell,1/4)\). The product basis changes take \(O(n^2\operatorname{poly}(\ell))\) gates.

In characteristic two, use the one-digit basis \(e_{01},e_{00},e_{11},e_{01}+e_{10}\) with weights \(-1,0,0,1\). The operator \(K\) sends the first basis vector to the last and kills the rest, while \(P=I+K\) never lowers weight. Take \(S\) to be the sum of coefficients through \(D\). Evaluation at \(t=1\) gives \(I\), so every nonzero entry of \(R=I-S\) raises weight by at least \(2(D+1)\). Let \(U\) project onto weights at least \(D\) and \(V\) onto weights at most \(-D\). Then \(R=UR+(I-U)RV\). Each summand factors through a space of dimension at most \(4^\ell\exp(-\alpha^2\ell/2+o(\ell))\), by Hoeffding's bound~\cite[Theorem~2]{hoeffding} for the independent uniform basis weights in \([-1,1]\). These factors need only coordinate projections, basis changes, and applications of \(I-S\).

For any \(\delta<\alpha^2/(2\ln2)\) the exceptional dimension is at most \(n^{2-\delta}\) for large \(\ell\). Keeping \(z\) makes the first rectangle group of size \(n^{3-\delta}\). For initial dimensions \((x,y,z')\), the three exponent bounds are \(x+y\) for preparation, \(h+H_2(\alpha)+2\alpha\) for the partial swaps, and \(\omega(x-\delta,y,z')\) for each low-rank term. Choose \((\alpha,\delta)=(1/25,1/1000)\) for \(k=7\) and \((1/80,1/10000)\) for \(k=8\). The respective rectangle exponents are \(15992001/2999000\) and \(1763916001/287990400\). All three bounds are strictly below \(16/3-1/1200\) and \(49/8-1/12000\), respectively. For the entropy comparison write \(\alpha=1/b\) and the available margin over \(h\) as \(p/c\). Raising to integer powers reduces the test to \((4b^b)^c<2^{bp}(b-1)^{(b-1)c}\).
\clearpage
\pdfbookmark[1]{References}{small-pattern-references}
\bibliographystyle{../linkedtitles}
\bibliography{../references}
\end{document}
